\begin{wbepi}{Jean-Jacques Rousseau (1762)}
When, among the happiest people in the world, bands of peasants are seen regulating 
affairs of State under an oak, and always acting wisely, can we help scorning the 
ingenious methods of other nations, which make themselves illustrious and wretched 
with so much art and mystery?
\end{wbepi}

I came to graduate school hoping to continue pursuing a question that weighed on 
my mind since high school -- why are school boards elected? In college I began 
by studying the differences between policy makers' imagined governance of school 
districts and their actual governance. In graduate school, I wanted to understand 
if school districts behaved democratically. And, after a long journey, I believe 
I have been able to provide a foundation from which to study this question. It 
has been a great pleasure to follow this line of thinking across more than a 
decade of study and to be satisfied with the destination. 

That journey of inquiry would not have been possible without the love, support, 
reassurance, and guidance of so many people. For intellectual inspiration I can 
thank Jeff Seward for encouraging me to pursue this work through graduate school 
and helping me find my way to UW-Madison. At UW-Madison I was able to and 
encouraged to pursue my own course of study, despite the relative disinterest 
in school board elections in the wider discipline. The faculty were always 
receptive to questions about how theories of political science might apply at this 
level of study, and I was lucky to spend time among scholars who emphasized 
proper measurement, robust analysis, and constructive criticism. 

Emblematic of this in so many ways is my advisor, John Witte. Not only was John 
supportive of my idea to pursue a topic outside of his interests, but he engaged 
with the subject of school boards with his trademark passion and enthusiasm. He 
introduced me to a number of scholars I could share ideas with and get feedback 
from. This outside expertise and enthusiastic support was critical when I made 
the decision to start work full-time at the Wisconsin Department of Public 
Instruction while just beginning to write my dissertation proposal. Much could 
be said about this decision, but the important part is that John assured me as 
I made it that he would support me however I needed in order to graduate, and 
he was true to his word. 

Working at DPI while completing the dissertation provided me with access to a 
wealth of expertise about Wisconsin school districts and a deeper understanding 
of relations between school boards, their district administrators, and the state 
and federal governments. 


The list of people who have made this dissertation possible is quite long. The 
dissertation has benefited from the help of numerous public officials in the 
state of Wisconson who generously lent their time and expertise to help better 
my understanding of the complex electoral process. Just a few of these individuals 
include John Ashley and Dan Rossmiller of the Wisconsin Association of School 
Boards, Molly Kuroda from the Wisconsin Department of Public Instruction, 
Jean Gottwald and Lori Stottler, county clerks, Vanessa Schwartz for her 
research assistance, colleagues at the Wisconsin DPI, and of course the staff 
at hundreds of Wisconsin school districts who helped find and transfer school 
board election records for the purposes of this study. 

My colleagues at the Department of Public Instruction.

The guidance and scholarship of other political scientists.

Friends and family. 

Last but not least, the patience of my wife - Hannah Miller - who supported me 
completing this project while working full time. 

Financial Assistance. 
The research reported here was supported by the Institute of Education Sciences, U.S. Department of Education, through Award # R305B090009 to the University of Wisconsin-Madison.  The opinions expressed are those of the authors and do not represent views of the U.S. Department of Education. 
Additional funding was provided through the University of Wisconsin-Madison 
Political Science department's summer initiatives funding for graduate students. 

